\documentclass{article}
\usepackage{dcase2025_techrep}
\usepackage{booktabs}
\usepackage{url}

\title{Audio-Dependency-Aware Post-Training of Qwen3-Omni for DCASE 2026 Task 5}
\name{Hyeonuk Nam}
\address{Independent researcher\\
Seoul, South Korea\\
frednam@kaist.ac.kr}

\begin{document}
\maketitle

\begin{abstract}
This technical report describes our system for DCASE 2026 Task 5, Audio-Dependent Question Answering.
We build on Qwen3-Omni-30B-A3B-Instruct~\cite{qwen3omni2025,qwen3omnihf}, an open-source large audio-language model, and focus on reducing shortcut reliance on the textual question prior.
Our strongest train-only system uses LoRA supervised fine-tuning on audio-dependent training items, a small number of empty-audio negative examples, and a text-only response normalizer for parse-failed model generations.
On the official development set, the best train-only system reaches 67.27\% accuracy after response normalization, compared with 65.90\% for our local Qwen3-Omni baseline.
Final submissions additionally include train+development variants and a two-model ensemble.
\end{abstract}

\section{Introduction}
DCASE 2026 Task 5 evaluates audio-dependent multiple-choice question answering on ADQA-Bench~\cite{dcase2026task5}.
The task differs from generic audio captioning or audio tagging because the model must answer a specific question from the audio, while many choices remain plausible from language priors alone.
In preliminary experiments, we observed that strong open audio-language models often answer a large fraction of samples correctly even when the audio is removed.
This motivated an audio-dependency-aware training recipe inspired by audio contribution analysis in AudioMCQ~\cite{he2026audiomcq}: we explicitly identify training samples for which the audio appears necessary, then fine-tune primarily on those items.

\section{System Overview}
All systems use Qwen3-Omni-30B-A3B-Instruct~\cite{qwen3omni2025,qwen3omnihf}.
We fine-tune using LoRA~\cite{hu2021lora} with rank 4, alpha 8, dropout 0.05, and trainable query and value projections only.
The audio tower and multimodal projector are frozen.
Inference uses greedy decoding with a prompt that asks the model to output only the exact option text.

For selected submissions, raw generations are post-processed in two stages.
First, deterministic parsing checks exact option-text matches and simple letter or prefix patterns.
Second, parse-failed responses are passed to the base Qwen3-Omni model as a text-only multiple-choice response normalizer.
This normalizer sees only the candidate choices and the raw generation; it does not receive the audio or the ground-truth label.
We count this normalizer as an additional 30B model component in the metadata.

\section{Diagnostic Audio-Dependency Analysis}
\label{sec:diagnostics}
Before fine-tuning, we ran diagnostic inference on train and development samples under multiple conditions: normal audio-question pairs, empty-audio questions, and shuffled-audio questions.
The goal was not to create a perfect taxonomy, but to separate samples likely solved by text priors from samples where the model depends on the audio.
Table~\ref{tab:buckets} shows the resulting bucket distribution.

\begin{table}[t]
\centering
\caption{Audio-dependency buckets from Qwen3-Omni diagnostic runs. Train buckets use the full train split; development buckets use the official development set.}
\label{tab:buckets}
\begin{tabular}{lrr}
\toprule
Bucket & Train & Dev \\
\midrule
Easy text-prior variants & 10648 & 402 \\
Strong audio-dependent & 4738 & 374 \\
Audio-helped / shuffle-leak & 1712 & 283 \\
Hard audio-dependent candidate & 1312 & 318 \\
Misleading audio or prior-only & 793 & 95 \\
Wrong normal but shuffle correct & 277 & 135 \\
\bottomrule
\end{tabular}
\end{table}

The train split contains 4,738 strong audio-dependent items, approximately 24.3\% of the full train set.
The development set contains 374 strong audio-dependent items, approximately 23.3\%.
High-volume categories include speaker identity, speech content, temporal reasoning, speech paralinguistics, music, and sound events.
This analysis also revealed that many development samples are highly prior-driven; therefore, optimizing only for the aggregate score can favor models that preserve text-prior behavior rather than models that improve audio grounding.

\section{Fine-Tuning Data Construction}
The main supervised fine-tuning set contains only strong audio-dependent training items.
We also tested broader subsets that include hard audio-dependent candidates and shuffle-leak items, but these did not improve development accuracy.
To discourage unsupported guessing, we add empty-audio negative examples.
For these examples, the audio input is removed and the target response is ``Cannot be determined from the audio.''
We sweep the negative-example ratio and the checkpoint step, using the development set for model selection.

\section{Development Results}
Table~\ref{tab:main_results} summarizes the main train-only systems.
The judge column applies the Qwen3-Omni text-only response normalizer only to parse-failed predictions.

\begin{table}[t]
\centering
\caption{Main train-only development results.}
\label{tab:main_results}
\begin{tabular}{lrrr}
\toprule
System & Step & Strict & Judge \\
\midrule
Qwen3-Omni baseline & -- & 65.90 & -- \\
Strong only & 3000 & 64.97 & 66.21 \\
Strong + empty 2.5\% & 3000 & 65.77 & 67.02 \\
Strong + empty 5\% & 1000 & 64.78 & 66.09 \\
Strong + empty 5\% & 2000 & 66.02 & 67.27 \\
Strong + empty 5\% & 3000 & 65.28 & 66.52 \\
Strong + empty 7.5\% & 2500 & 65.34 & 66.52 \\
Strong + empty 10\% & 1000 & 64.90 & 66.15 \\
Strong + empty 20\% & 3000 & 65.09 & 66.21 \\
Strong + empty+shuffle 10\% & 1000 & 64.97 & 66.27 \\
\bottomrule
\end{tabular}
\end{table}

The best train-only setting is the 5\% empty-audio recipe at 2000 steps.
The 2.5\% recipe at 3000 steps is the second strongest setting after response normalization.
Higher empty ratios and combined empty+shuffle negatives did not outperform the best two settings.
Letter-only and explicit chain-of-thought inference prompts were also worse than exact-option generation: for the best checkpoint, letter-only reached 60.17\% after normalization and CoT-then-letter reached 65.34\%, both below the default prompt.

\section{Additional Ablations and Failed Directions}
\label{sec:failed}
We evaluated several alternative improvement directions; none replaced the simple strong-audio SFT recipe.
Table~\ref{tab:failed} lists the main negative results.

\begin{table}[t]
\centering
\caption{Representative failed or deprioritized experiments on the development set.}
\label{tab:failed}
\begin{tabular}{lrr}
\toprule
Experiment & Strict & Judge \\
\midrule
FunAudioChat base & 54.26 & -- \\
AudioFlamingo3 rerun & 54.32 & -- \\
Answer-only Qwen SFT, 9k & 62.85 & -- \\
Silent CoT SFT, 3k & 61.05 & -- \\
Explicit CoT SFT, 3k & 58.49 & -- \\
Category-balanced v3 & 64.09 & -- \\
Catbalmix10, ckpt2000 & 65.21 & -- \\
SIMPO continuation, 500 & 65.65 & 66.96 \\
DPO continuation, 500 & 65.15 & 66.46 \\
SFT-init GRPO, 100 & 65.09 & 66.21 \\
SFT-init DAPO-lite, 100 & 65.21 & 66.40 \\
Empty-aware GRPO, 100 & 65.09 & 66.46 \\
Empty-aware DAPO-lite, 100 & 64.72 & 65.90 \\
\bottomrule
\end{tabular}
\end{table}

The early answer-only Qwen3 fine-tuning runs degraded the base model, even though the base model was already the strongest single model among the tested backbones.
Category-level model routing provided only a small oracle gain: choosing the best model per rule-based category would reach 66.83\%, compared with 65.90\% for Qwen3-Omni base.
The much higher per-sample oracle of 90.85\% indicates complementary errors, but coarse category routing alone is insufficient.
Preference/RL continuation also failed to exceed the 5\% empty-audio SFT checkpoint.
ORPO collapsed into invalid outputs in our run, and KTO was not used because the audio metadata path was dropped by the training pipeline.

\section{Error Analysis}
Comparing the best train-only system with the Qwen3-Omni baseline shows a clear trade-off.
After response normalization, the fine-tuned model fixes 130 baseline errors but introduces 108 regressions.
Many gains are in speech/voice/phonetics, music, counting, and sound-event questions.
Regressions are also concentrated in speech/voice/phonetics and music, suggesting that fine-tuning improves audio dependence but can damage some prior knowledge or fine-grained perception.
In the bucket analysis, the fine-tuned model improves hard-all-wrong and misleading-prior buckets, but loses accuracy on text-prior-easy and shuffle-leak buckets.
This supports our final choice to include one train-only system and two train+development variants rather than relying on a single heavily specialized model.

\section{Submitted Systems}
We prepare four candidate submissions.
System 1 is the train-only strong-audio SFT checkpoint with 5\% empty negatives at 2000 steps and base Qwen3 response normalization.
System 2 retrains the 5\% recipe on the union of train and development strong audio-dependent items.
System 3 uses the train+development recipe with 2.5\% empty negatives and a 3000-step checkpoint.
System 4 ensembles Systems 2 and 3; when the two systems disagree, System 2 is used as the tie breaker.
Because the development labels are used in Systems 2--4, no held-out development score is reported for those final variants.

\section{Conclusion}
Our final system is a simple Qwen3-Omni LoRA pipeline guided by audio-dependency diagnostics.
The most effective intervention is not a more complex backbone or RL objective, but careful removal of easy text-prior samples and a small amount of empty-audio negative training.
Although the improvements are modest, the experiments suggest that explicit audio-dependency filtering is a practical way to adapt large audio-language models for ADQA-style evaluation.

\bibliographystyle{IEEEtran}
\bibliography{references}

\end{document}
